\documentclass[11pt,letterpaper]{article}
\usepackage[T1]{fontenc}
\usepackage[utf8]{inputenc}
\usepackage{lmodern}
\usepackage[margin=1in]{geometry}
\usepackage{amsmath,amssymb,amsthm,mathtools}
\usepackage{booktabs,array,longtable}
\usepackage{xcolor}
\usepackage{listings}
\usepackage{microtype}
\usepackage{enumitem}
\usepackage[colorlinks=true,linkcolor=blue!45!black,urlcolor=blue!45!black,citecolor=blue!45!black]{hyperref}
\usepackage{bookmark}
\usepackage{fancyhdr}
\setlength{\headheight}{14pt}
\pagestyle{fancy}
\fancyhf{}
\fancyhead[L]{\small Concise iPRA: implementation and closure}
\fancyhead[R]{\small Research note}
\fancyfoot[C]{\thepage}
\setlength{\parskip}{4pt}
\setlength{\parindent}{0pt}
\setlist{itemsep=3pt,topsep=4pt}
\emergencystretch=2em
\newtheorem{theorem}{Theorem}[section]
\newtheorem{lemma}[theorem]{Lemma}
\newtheorem{proposition}[theorem]{Proposition}
\newtheorem{corollary}[theorem]{Corollary}
\theoremstyle{definition}
\newtheorem{definition}[theorem]{Definition}
\theoremstyle{remark}
\newtheorem{remark}[theorem]{Remark}
\newcommand{\file}[1]{\expandafter\nolinkurl\expandafter{\detokenize{#1}}}
\lstset{language=Python,basicstyle=\ttfamily\footnotesize,
  columns=fullflexible,keepspaces=true,showstringspaces=false,
  breaklines=true,breakatwhitespace=false,frame=single,
  rulecolor=\color{black!20},xleftmargin=4pt,xrightmargin=4pt}
\hypersetup{pdftitle={A Smaller Implementation and a Deterministic Closure Presentation for iPRA},pdfauthor={Research note prepared for Vladimir Reshetnikov}}
\title{\bfseries A Smaller Implementation and a Deterministic\\[4pt]
Closure Presentation for iPRA}
\author{Research note prepared for Vladimir Reshetnikov}
\date{10 November 2026}
\begin{document}
\maketitle
\begin{abstract}
The requested strictly simpler mathematical axiomatization has not been
established. The defensible results are narrower. First, a local change to
two complete Python proof checkers for the existing native scalar calculus
reduces the inspected byte minimum from 529 to 528 and the inspected AST
minimum from 289 to 288, in different programs. The accepted certificates
are unchanged in the ideal resource-unbounded reading. Token counts and
gzip sizes increase in the matched comparisons.

Second, a different native presentation has a deterministic introduction
rule: literal-equation closures express progress, while compound closures
retain the original conditional-universal meaning. Compound padding gives
a proof translation from the existing core, and explicit formula composites
establish a whole-domain bi-interpretation with iPRA. Its more elaborate
elimination rule has not been shown to reduce total complexity. Finally, an
exact zero-set theorem explains why fixing every scalar iterator's initial
state to zero cannot preserve the existing term compiler. Complete rules,
written proofs, executable checks, provenance, and the inherited arithmetic
bridge are supplied.
\end{abstract}
\tableofcontents
\newpage

\section{Outcome, comparison target, and scope}

The comparison target is the current \file{ConciseIPRA} project at
\href{https://github.com/VladimirReshetnikov/ProveIt/tree/8662bc0a00f4e77e482f7d141d7a4f919f70f9aa/Logic/PrimitiveRecursiveArithmetic/ConciseIPRA}{commit \texttt{8662bc0a00f4}}.
The full commit is recorded in the package provenance. Its README describes
the native scalar core and 22 filed research packages, with complete
standalone-checker measurements and an ongoing Lean development
\cite{current}. All comparisons in this note use that fixed snapshot;
they do not make a claim about later changes or a global minimum.

\begin{center}
\begin{tabular}{p{0.25\linewidth}p{0.65\linewidth}}
\toprule
Item & Established status\\
\midrule
Existing-core implementation & Smaller by one byte or one AST node in
separate complete programs; no change to the calculus.\\
Deterministic mixed closure & A different presentation, with written
proofs of both interpretations and their composites. No smaller checker
for it has been demonstrated.\\
Zero-start iterator deletion & A same-domain term-definability obstruction;
not an impossibility theorem for arbitrary interpretations.\\
Requested axiomatic improvement & Still unresolved by this investigation.\\
\bottomrule
\end{tabular}
\end{center}

The distinction between these outcomes matters. Merely displaying one
introduction schema can move its cost into another rule. Likewise, a
one-byte implementation change does not eliminate an axiom or reduce
arithmetical strength. This note reports the exact improvement it can
justify and leaves the stronger requested conclusion open.

\subsection{The common arithmetic and logical baseline}

The current numerical terms are
\[
 t ::= x\mid0\mid St\mid I(n,a;u.h),
\]
where only the update body $h$ binds $u$. There are two computation schemas:
\[
 I(0,a;u.h)=a,
 \qquad I(Sn,a;u.h)=h[I(n,a;u.h)/u].
\]
The native formula vocabulary uses equality and the binders $Q,D$, with
ordinary readings
\[
 Q_yB(y)=B(0)\to\forall yB(y),\qquad
 D_{x,y}B(x,y)=\exists x\,Q_yB(x,y).
\]
The designated contradiction is $0=S0$. The core has hypotheses,
Leibniz transport through arbitrary native formula templates, conditional
universal elimination, fresh-witness elimination, and the two introduction
patterns
\[
 \frac{\Gamma,B(t,0)\vdash B(t,Sz)}{\Gamma\vdash D_{x,y}B(x,y)},
 \qquad
 \frac{\Gamma,E(t,z)\vdash E(t,Sz)}{\Gamma\vdash D_{x,y}E(x,y)},
\]
where the latter requires a literal equation and $z$ is fresh. The full
inventory and all eigenvariable conditions are reproduced in
Appendix~\ref{sec:choice-introduction}; that equivalent presentation treats
$Q$ assumptions as reusable conditional clauses. This is inherited
mathematics, not a new reduction.

The ordinary realization $T_{\mathrm{eq}}$ has intuitionistic first-order
equality logic, the same iterator schemas, numerical coverage,
$(0=S0)\to\bot$, and contextual induction for literal equations. The
inherited bridge constructs predecessor, equality reflection,
quantifier-free induction, pairing, and a compiler for each fixed
primitive-recursive definition. It proves both internal numerical
composites, rather than only agreement in the standard model. The relevant
proofs are reproduced in the appendices and retain their source
attribution. In particular, this note does not replace full first-order
iPRA with its equational or closed-numeral fragment.

\section{The smaller complete implementations}

\subsection{Measurements and the unchanged public contract}

The complete module is counted: JSON input, substitution/opening,
contexts, rule checking, recursion, and the final verdict. The common
infrastructure is CPython and its standard JSON library, exactly as in the
comparison sources. A run accepts precisely when its exit status is zero
and its standard output is the exact line \file{ACCEPT}. Uncaught checking
or parsing exceptions reject. Neither a hidden checker import nor an
external executable implements any proof rule.

Measurements use CPython 3.12.14. Parse the whole file with \file{ast.parse},
normalize with \file{ast.unparse}, and append one newline. Significant
tokens exclude comments, indentation, line breaks, and stream markers;
AST nodes are all nodes visited by \file{ast.walk}. Gzip uses compression
level 9 and \file{mtime=0}; it is an informational coordinate. The two
candidate source files are already in normalized form.

\begin{center}
\begin{tabular}{lrrrr}
\toprule
Complete program & Tokens & AST nodes & Bytes & Gzip 9\\
\midrule
Reference AST profile &263&289&613&349\\
New AST profile       &264&\textbf{288}&608&350\\
Reference byte profile&266&310&529&333\\
New byte profile      &267&309&\textbf{528}&336\\
\bottomrule
\end{tabular}
\end{center}

The references are respectively
\file{coverage_and_kernel/code/checker_ast.py} and
\file{neutral_scope/code/checker.py} \cite{coverage,scope}.
The candidate modules are \file{code/checker_ast.py} and
\file{code/checker_bytes.py}. The old 261-token coordinate remains
unbeaten. No one delivered program has all three coordinate minima.
The gzip figures increase in both matched comparisons. These are
measurements of particular files, not lower bounds on possible checkers.

\subsection{The exact transformation}

An instruction checks a major premise, optionally changes the current goal
by transport or closure specialization, computes an available hypothesis,
and opens the next scope. The references execute the major check as a
standalone statement:
\begin{lstlisting}
c(m, g, k)
# Pure transport/specialization precheck and possible goal replacement.
h = m * 3 or ...
\end{lstlisting}
The new versions execute that pure precheck first and place the same major
check at the start of the existing expression:
\begin{lstlisting}
# The same pure precheck and possible goal replacement.
h = c(m, g, k) or m * 3 or ...
\end{lstlisting}
Every successful invocation of \file{c} returns \file{None}, so the inserted
operand does not change the value of \file{h}. The major proof is still
checked before that value is constructed. Crucially, the later scope-floor
decrement is not moved. Witness opening retains its original order relative
to fresh-variable allocation.

The idea of prevalidating an elimination template is already present in the
repository's \file{prevalidated_elimination} package \cite{prevalidate}.
The contribution here is its small composition with these lexical-scope
implementations. It is not claimed as a new proof-theoretic idea.

\begin{theorem}[Equality of accepted certificates]
For each of the two matched reference/candidate pairs, the same finite
JSON certificates reach the accepting verdict in idealized Python
execution, ignoring physical time, memory, and stack limits.
\end{theorem}
\begin{proof}
All parsed data are ordinary JSON values and built-in containers; there
are no user-supplied operator overloads, callbacks, or threads. The moved
precheck reads the instruction, current goal, and local opening function.
It neither consumes the certificate stream nor mutates the data, retained
context, or scope floor. Its only caller-local change is replacement of
the current goal.

The major check receives exactly the same arguments $m,g,k$ in either
program. It has its own local floor and local context binding, and it
does not refer to the caller's goal. The byte version extends contexts
by list concatenation; the AST version uses tuple concatenation. Neither
changes its caller's context. The relevant shared effect is consumption
of the certificate stream, which the pure precheck never accesses.

After a successful major check, its return value is \file{None}. For the
recursive version this follows from its bare or implicit return; for the
loop version it follows from falling out of the loop. Thus the leading
operand in the existing \file{or} expression is false and evaluation
continues with exactly the old expression. The following state, goal,
scope floor, and remaining input stream are the same.

Induct on a finite successful checking tree. Reorder the precheck and
major call at its root, apply the induction hypothesis to recursive major
calls and continuations, and obtain the same accepting execution. The
argument works in both directions. If either local computation raises
an error, it is not swallowed by this transformation and the final verdict
cannot be reached. The first exception or rejection latency may change,
which is immaterial to the stated acceptance contract.
\end{proof}

The theorem transfers the existing mathematical soundness and completeness
claim to the candidates; it does not independently reprove every invariant
of the inherited permissive certificate decoder. In particular, the AST
profile retains its existing empty-string and empty-object truth aliases.
These are recorded in the tests, rather than silently changing the input
language during comparison.

\subsection{Executable evidence and provenance}

The delivered \file{verify.py} checks 49 inherited certificates and 14
boundary/alias cases against all four complete programs, under both
ordinary and optimized Python: 504 command-line runs. The outcome was
zero mismatches. Six additional mutation-control runs deliberately
introduced three known unsound changes: a shared mutable context,
an incorrect induction selector, and allocating a fresh parameter before
opening its witness. Each corresponding false-proof control was accepted
by the deliberately broken program, while the actual programs rejected it.
The complete per-run ledger is \file{tests/results.json}.

The inherited certificates exercise arithmetic computation, equality
transport, literal-equation induction, arbitrary-term $Q$ specialization,
compound cases, and nonvacuous existential witnesses. Negative cases cover
failed majors, false equations, forbidden compound induction, escaping
parameters, context leaks, and numeric aliases. Finite tests are evidence
for implementation fidelity; the universal argument is the source theorem
above and the inherited calculus proof.

An exact normalized-AST screen covered 4,532 Python file occurrences,
representing 549 distinct raw sources, in the pinned Research tree.
Neither candidate had an exact match, and no file failed to parse.
This establishes only absence of an exact normalized program in that
screened corpus. It does not establish conceptual novelty; the prior
prevalidation work has already been acknowledged.

\subsection{Complete byte-profile source}

The following listing is the complete 528-byte normalized module. Line
wrapping is typographical and is not part of the measured file.
% BYTE_LISTING_INSERTION

\section{How to use the package and what remains open}

\begin{lstlisting}
python3 measure.py
python3 verify.py
python3 code/checker_bytes.py < tests/example_accept.json
pdflatex -interaction=nonstopmode -halt-on-error report.tex
pdflatex -interaction=nonstopmode -halt-on-error report.tex
\end{lstlisting}

The first command reproduces the complete-module metrics. The second
recreates the executable validation ledger. The third runs an actual
certificate through the new byte profile. The report's TeX source is
self-contained: its mathematical sections and displayed checker source
are embedded, so it does not require fetching a bibliography, repository,
or other document. TeX packages and a Python installation are the ordinary
build dependencies. The ZIP includes the separate section sources for
convenient review as well.

\paragraph{Further research questions.}
\begin{enumerate}
\item Can the matrix-sensitive closure in the next section be implemented
with less total checking work than the existing mode test? Its source
and target instantiations, iterator construction, root test, and public
driver must all be included in the comparison.
\item Can a graph interpretation recover arbitrary initialization from
a smaller numerical syntax while internally proving totality,
functionality, induction transfer, and both interpretation composites?
The zero-set obstruction below excludes a direct term compiler, but
does not settle this more general question permitted by the request.
\item Can the mixed-closure translation and the same-category congruence
lemma be formalized in the repository's existing Lean framework? A
formalization should expose the syntax-sensitive meaning of closure
and forbid cross-category congruence explicitly.
\item Is there a mathematical lower bound for a carefully fixed class
of standalone proof presentations? A census of existing source files
is not such a lower bound, and unrestricted ``simplest'' remains
dependent on the chosen presentation language and size convention.
\end{enumerate}

The next two sections contain the new mathematical alternatives and
obstruction. The appendices reproduce the established native calculus
and arithmetic bridge with attribution. The new arguments are written
proofs reviewed for circularity and scope conditions; no new Lean or
Rocq verification is claimed.

